% Generated draft: requires mathematical and source review.
\documentclass[11pt]{article}
\usepackage[margin=25mm]{geometry}
\usepackage{fontspec}
\setmainfont{Latin Modern Roman}
\usepackage{amsmath,amssymb,mathtools,unicode-math}
\setmathfont{Latin Modern Math}
\usepackage{enumitem,needspace,xurl,hyperref}
\setlength{\parindent}{0pt}
\setlength{\parskip}{5pt}
\setlength{\emergencystretch}{3em}
\newcommand{\N}{\mathbb{N}}
\newcommand{\Z}{\mathbb{Z}}
\newcommand{\Q}{\mathbb{Q}}
\newcommand{\R}{\mathbb{R}}
\newcommand{\C}{\mathbb{C}}
\newcommand{\F}{\mathbb{F}}
\newcommand{\A}{\mathbb{A}}
\newcommand{\PP}{\mathbb{P}}
\newcommand{\E}{\mathbb{E}}
\newcommand{\eps}{\varepsilon}
\newcommand{\dd}{\,\mathrm{d}}
\begin{document}
\section*{UnsolvedMath ID: 800877}
\textbf{Source problem code:} catalogue entry
\subsection{Definitions and mathematical statement}
Let \(X=\C\mathbb{P}^{n}\), with \(n\geq 1\), and let \(\omega\) be its Fubini--Study K\"ahler form. An \(\omega\)-plurisubharmonic function is an upper semicontinuous, locally integrable function \(u:X\to[-\infty,\infty)\) for which \(\omega+dd^{c}u\) is a positive current. Here \(dd^{c}\) and the scale of \(\omega\) are to follow the source conventions, which the supplied excerpt does not specify.

The global domain of definition \(DMA(X,\omega)\) consists of those \(\omega\)-plurisubharmonic functions \(u\) for which every decreasing sequence of bounded \(\omega\)-plurisubharmonic functions \(u_j\downarrow u\) has Bedford--Taylor measures \((\omega+dd^{c}u_j)^n\) converging weakly to a positive measure independent of the approximating sequence. Denote that measure by \((\omega+dd^{c}u)^n\). Write \(\delta_a\) for the unit Dirac measure at \(a\), and \(\nu(u,a)\) for the Lelong number at \(a\) of a local plurisubharmonic potential for \(\omega+dd^{c}u\); adding a smooth local potential does not affect this number.

\textbf{Question (attributed in the source to Coman--Guedj).} Can one find a point \(a\in X\) and a function \(G_a\in DMA(X,\omega)\) such that
\[
 (\omega+dd^{c}G_a)^n
 =\left(\int_X\omega^n\right)\delta_a,
 \qquad
 \nu(G_a,a)\in\R\setminus\Q?
\]
Thus the entire Monge--Amp\`ere mass must be concentrated at the single point \(a\). The volume factor retains the unspecified normalization of the Fubini--Study form. The excerpt does not state whether examples are sought in some dimension, in every dimension, or for a particular fixed dimension; \(n\) is therefore left as a dimension parameter. The normalization of Lelong numbers must also be checked against the original source before interpreting the numerical irrationality condition.
\subsection{Short English statement}
The question attributed to Coman--Guedj asks for an \(\omega\)-plurisubharmonic function on complex projective space, in the global domain of definition of the Monge--Amp\`ere operator, whose entire Monge--Amp\`ere measure is supported at one point and whose Lelong number there is irrational. The dimension quantifier and numerical normalizations require source confirmation. The registry labels the problem open; that label is supplied metadata, not an independently verified assessment.
\subsection{Sources}
\begin{itemize}
\item Local UnsolvedMath registry: permanent ID \(800877\), exact source problem code \texttt{catalogue entry}, title \emph{Open Problems in Pluripotential Theory --- Question 8}. The supplied record lists Dinew, Guedj, Zeriahi as proposers and gives the year 2015.
\item Dinew, Guedj, Zeriahi, \emph{Open Problems in Pluripotential Theory} (2015), Question 8: \url{https://arxiv.org/abs/1511.00705}. The question itself carries the attribution \textnormal{[Coman--Guedj]} in the supplied excerpt.
\item Supplied provenance: accessed 2026-10-16; extraction method \texttt{arxiv\_latex}; upstream revision \texttt{d371ee917ac9647fbb06cd6fc529deec6f3dda82}. These are supplied provenance details; no source access or literature verification was performed for this draft.
\end{itemize}
\end{document}
